\chapter{Results and Evaluation}

Ninety pages from three thesis documents were converted by each of the four
strategies and scored against hand-written references. The results show a
consistent pattern across all three documents and address all five research
questions.

\section{Experiment Setup}

The experiment evaluates four extraction strategies against a
hand-written ground truth of 90 Markdown files, 30 per thesis document,
as described in Section~4.2. The ground truth was written manually to
avoid circularity: an automatically generated reference would embed the
same errors as the system under evaluation.

The three documents in the corpus are a German computer science
bachelor thesis on neural networks (referred to here as the neural
networks thesis), an English-language bachelor thesis on sentiment analysis
(the sentiment analysis thesis), and another English-language bachelor
thesis describing a human-in-the-loop system for research paper
generation (the HITL thesis). Together they cover different writing
styles, content densities, and proportions of code, formula, and
table content, which makes the per-page results more representative
than a single document would allow.

All four strategies (text-only, image-only, hybrid, and adaptive) were
applied to every page of every document. This produced
$4 \times 3 \times 30 = 360$ evaluated conversions with zero errors.
For each conversion, the pipeline output was compared to the
corresponding reference file using the seven metrics described in
Section~4.3: text similarity, heading structure, list structure,
table structure, code block score, paragraph structure, and word
overlap. Processing time was also recorded per page.
Results are reported as mean scores per strategy across the full
90-page corpus.

\section{Overall Strategy Performance (RQ1, RQ2)}

Table~\ref{tab:results_overview} shows the mean scores for all four
strategies across the full corpus.

\begin{table}[ht]
\centering
\caption{Mean metric scores per strategy across 90 pages (3 documents)}
\label{tab:results_overview}
\renewcommand{\arraystretch}{1.4}
\small\setlength{\tabcolsep}{4pt}
\begin{tabularx}{\textwidth}{lXXXXXXX}
\toprule
\textbf{Strategy} & \textbf{Text sim.} & \textbf{Heading} & \textbf{Lists} & \textbf{Tables} & \textbf{Code} & \textbf{Paras.} & \textbf{W. overlap} \\
\midrule
Text-only & 0.893 & \textbf{0.808} & 0.904 & 0.933 & \textbf{1.000} & \textbf{0.815} & 0.804 \\
Adaptive  & \textbf{0.894} & 0.803 & \textbf{0.919} & \textbf{0.978} & 0.994 & 0.808 & \textbf{0.854} \\
Image-only & 0.839 & 0.685 & 0.911 & 0.944 & 0.972 & 0.782 & 0.838 \\
Hybrid    & 0.813 & 0.523 & 0.875 & 0.944 & 0.946 & 0.774 & 0.835 \\
\bottomrule
\end{tabularx}
\end{table}

\begin{figure}[ht]
    \centering
    \includegraphics[width=\textwidth]{performance.png}
    \caption{Mean metric scores per strategy across 90 pages}
    \label{fig:performance}
\end{figure}

Figure~\ref{fig:performance} shows the mean scores per strategy across all seven metrics.
Text-only and adaptive are close on most metrics, close enough that either
could be called the stronger baseline depending on which one matters most.
Text-only leads narrowly on heading structure (0.808), code blocks (1.000),
and paragraph structure (0.815). Adaptive leads on list structure (0.919)
and word overlap (0.854), and the two are effectively tied on text
similarity (0.894 versus 0.893). Text-only still costs a fraction of a
second per page and never touches the LLM; its strength comes from the same
source as before: thesis documents are mostly continuous prose, and
PyMuPDF's text extraction is very accurate on that kind of content.

Adaptive's clearest advantage is table structure. Its score of 0.978 is
well ahead of every fixed strategy, including text-only at 0.933 and
image-only and hybrid at 0.944. Pages classified as TABLE or MIXED are
routed to the vision model, which reads a table's visual grid directly
instead of relying on the PDF's text layer, and that shows up as the widest
gap anywhere in the results table.

Image-only and hybrid both trail on most metrics, but not by as much as an
uncorrected reading of the raw model output would suggest.
Section~\ref{sec:output_quality} covers this in more detail: image-only's
paragraph and code block scores (0.782 and 0.972) are now close to the other
three strategies, and heading structure (0.685) is its real remaining
weakness. Hybrid's clearest weak point is also heading structure (0.523),
the lowest score anywhere in the table.

These results partially answer RQ1 and RQ2. Adaptive does not beat
text-only outright, but it comes within noise of every metric text-only
leads and pulls well ahead on table structure and word overlap, the two
metrics most sensitive to visual layout rather than plain prose. Section~\ref{sec:efficiency}
shows adaptive reaches this at a fraction of the cost of the two other
vision-based strategies.

\section{Efficiency (RQ3)}
\label{sec:efficiency}

Table~\ref{tab:efficiency} shows the mean processing time per page for
each strategy.

\begin{table}[ht]
\centering
\caption{Mean processing time per page by strategy}
\label{tab:efficiency}
\renewcommand{\arraystretch}{1.4}
\begin{tabular}{lrr}
\toprule
\textbf{Strategy} & \textbf{Avg time (ms)} & \textbf{Relative to text-only} \\
\midrule
Text-only  &   153.7 & $1\times$ \\
Adaptive   &   704.5 & $4.6\times$ \\
Image-only & 3{,}250.2 & $21.2\times$ \\
Hybrid     & 3{,}170.2 & $20.6\times$ \\
\bottomrule
\end{tabular}
\end{table}

\begin{figure}[ht]
    \centering
    \includegraphics[width=0.85\textwidth]{efficiency.png}
    \caption{Mean processing time per page by strategy (log scale)}
    \label{fig:efficiency}
\end{figure}

Figure~\ref{fig:efficiency} shows the mean processing time per page for each strategy on a log scale.
Text-only is by far the fastest strategy at 154\,ms per page on average.
It makes no LLM calls, so processing time is bounded by disk read and
PyMuPDF parsing speed. The adaptive strategy averages 705\,ms per page,
about 4.6 times slower than text-only but roughly 4.5 times faster than
either image-only or hybrid. The speedup comes directly
from the page classifier: TEXT pages skip the LLM entirely, and only pages
classified as IMAGE, TABLE, FORMULA, or MIXED go through the vision model.

Image-only and hybrid are the slowest strategies, at 3{,}250\,ms and
3{,}170\,ms respectively. Both send a full page image to the vision model
regardless of what the page contains, which costs far more than the
classifier-gated adaptive approach. Hybrid's extra text payload does not
add a meaningful further cost on top of that, and it still does not
translate into a consistent accuracy improvement.

Token usage tells the same story from a different angle.
Table~\ref{tab:tokens} shows the total tokens consumed across all 90 pages
and the average per LLM call.

\begin{table}[ht]
\centering
\caption{Token usage per strategy across 90 pages}
\label{tab:tokens}
\renewcommand{\arraystretch}{1.4}
\begin{tabular}{lrrrr}
\toprule
\textbf{Strategy} & \textbf{LLM calls} & \textbf{Avg tokens / call} & \textbf{Total tokens} \\
\midrule
Text-only  & 0  & ---     & 0       \\
Adaptive   & 23 & 3{,}669 & 84{,}397  \\
Image-only & 90 & 3{,}745 & 337{,}055 \\
Hybrid     & 90 & 4{,}163 & 374{,}661 \\
\bottomrule
\end{tabular}
\end{table}

The per-call token cost is similar for all three VLM strategies, ranging
from 3{,}669 to 4{,}163 tokens per call. The difference is in how many
calls are made. Image-only and hybrid call the LLM on every one of the
90 pages. Adaptive calls it on only 23 pages --- the pages the classifier
identified as IMAGE, FORMULA, TABLE, or MIXED. The other 67 pages are
TEXT pages that go through pymupdf4llm with no LLM call at all.
This reduces total token consumption by 75\% compared to image-only and
77\% compared to hybrid. Hybrid consumes the most tokens of any strategy
because it sends both the rendered page image and the extracted text in
every request, which increases the input length on every call.

These results confirm hypothesis H3: the adaptive system is significantly
cheaper than image-only or hybrid in both processing time and token usage,
while matching or beating them on most metrics. The efficiency gain grows
with document length --- a 300-page document would consume roughly
$300 \times 3{,}745 \approx 1.12$\,million tokens with image-only but only
around $281{,}000$ with adaptive, assuming a similar content distribution.
This is the same 75\% reduction observed on the 90-page corpus, scaled up.

\section{Per-Page Analysis and Consistency (RQ5)}
\label{sec:rq5_consistency}

The per-page breakdown shows that no single strategy wins every page.
Across the 90 evaluated pages, image-only produces the highest text
similarity on 19 pages, hybrid on 16, text-only on 11, and adaptive on 6.
On the remaining 38 pages, two or more strategies scored equally.
The variation reflects the content mix within each document.
Pages with complex table or figure layouts favour hybrid and image-only.
Dense text pages favour text-only or adaptive.

HITL page 9 shows the widest spread between strategies on a single page:
adaptive scores 0.968, text-only 0.664, while image-only and hybrid collapse
to 0.047 and 0.032. The reference file is a List of Tables, a short
front-matter list set as a pipe table. PyMuPDF's table finder classifies
this page as TABLE, and since pymupdf4llm already extracts a correct pipe
table from it, adaptive takes the text-only path described in
Section~\ref{sec:postprocessing_pipeline} instead of paying for a vision
call it does not need. Image-only and hybrid have no such option: both
render the page and ask the vision model to reconstruct the same list from
a picture of it, and the model reproduces almost none of it correctly.

An earlier version of this evaluation attributed adaptive's low score on
table-of-contents continuation pages to its design choice of leaving them
empty by default; correcting the evaluation to properly credit the
postprocessing pipeline showed this cost was smaller than originally
measured. HITL page 9 shows the opposite side of the same design: routing a
reliable page away from the vision model is a clear win here, not a cost.

The strategy rankings are close across all three documents. Adaptive leads
or matches text-only on table structure in every document: 1.000 versus
0.933 on the neural networks thesis, an exact tie at 0.933 on the HITL
thesis, and 1.000 versus 0.933 again on the sentiment analysis thesis. Text
similarity is close enough between the two that no consistent winner
emerges: adaptive is very slightly ahead on the neural networks thesis
(0.933 versus 0.928) and the HITL thesis (0.920 versus 0.919), while
text-only is very slightly ahead on the sentiment analysis thesis (0.833
versus 0.830). None of these gaps exceed half a percentage point. This
answers RQ5 more cleanly than the corpus-level averages alone: adaptive's
table structure advantage holds in every document, and the remaining
per-document variation in text similarity is small enough to read as noise
rather than a real reversal. A document with more formula or image pages
would widen the gap between adaptive and text-only because the classifier
would route more pages to the vision model.

\section{Output Quality Metrics (RQ4)}
\label{sec:output_quality}

The seven-metric framework still earns its keep even though no strategy
collapses the way an early, uncorrected reading of this evaluation once
suggested. Hybrid's heading structure score (0.523) sits nearly 0.3 below
its own text similarity score (0.813), a gap a single aggregate score would
average away entirely. A RAG system that chunks a document by heading level
would still misplace hybrid's output even though its overall similarity to
the reference looks reasonable on its own.

Heading structure remains the weakest metric relative to each strategy's
own other scores, though the picture is no longer uniform across
strategies. Text-only (0.808) and adaptive (0.803) both reach a respectable
score once postprocessing correctly reassigns heading levels from
font-size and section-number patterns. Image-only (0.685) and especially
hybrid (0.523) lag further behind, since both rely on the vision model to
infer heading depth from the page's visual layout rather than read it from
an explicit numbering pattern. This suggests heading level detection is
still the hardest structural task for any strategy that has to guess depth
from appearance rather than read it directly from the text.

Table structure is the metric where the strategies diverge most clearly.
Adaptive scores 0.978, ahead of every fixed strategy: text-only at 0.933,
and image-only and hybrid both at 0.944. On pages where text extraction
produces malformed or missing table markup, the vision model restores it,
and the classifier sends exactly those pages there while leaving simpler
tables in the faster text path (Section~\ref{sec:postprocessing_pipeline}).
This is the metric where per-page routing provides a clear benefit beyond
what either fixed strategy reaches alone.

Word overlap scores are close across all four strategies, ranging from
0.804 (text-only) to 0.854 (adaptive). This means all strategies recover
most of the correct vocabulary. The differences between strategies are in
structure and formatting, not in which words appear in the output.
